\documentclass[11pt]{article}
\usepackage[T1]{fontenc}
\usepackage{lmodern}
\usepackage[left=1in,right=1in,top=0.85in,bottom=0.85in]{geometry}
\usepackage{amsmath,amssymb,amsthm,mathtools}
\usepackage{microtype}
\usepackage[colorlinks=true,linkcolor=blue!50!black,urlcolor=blue!50!black,citecolor=blue!50!black]{hyperref}
\usepackage{xcolor}
\usepackage{enumitem}
\setlist{itemsep=3pt,topsep=5pt}
\setlength{\emergencystretch}{2em}
\newtheorem{theorem}{Theorem}
\newtheorem{lemma}[theorem]{Lemma}
\newtheorem{proposition}[theorem]{Proposition}
\theoremstyle{remark}
\newtheorem{remark}[theorem]{Remark}
\newcommand{\PA}{\mathsf{PA}}
\newcommand{\PAbin}{\mathsf{PA}_{\mathrm{bin}}}
\newcommand{\Prov}{\operatorname{Prov}}
\newcommand{\Prf}{\operatorname{Prf}}
\newcommand{\Len}{\operatorname{Len}}
\newcommand{\Chk}{\operatorname{Check}}
\newcommand{\num}[1]{\overline{#1}}
\newcommand{\gn}[1]{\ulcorner #1\urcorner}
\newcommand{\bits}{\operatorname{bits}}
\newcommand{\R}{\mathcal R}
\title{A slow-verifier separation for L\"ob's rule\\
\large An affirmative construction for the chosen-system clause of MAIS-O11(1)}
\author{Research note for Andrew Hagy\\\small AI-assisted mathematical draft; no Lean certificate}
\date{22 September 2026}
\begin{document}
\maketitle
\begin{abstract}
Assuming the consistency of Peano arithmetic, we construct a consistent,
recursively presented extension $S$ of PA, with binary numerals and nested,
priced abbreviations, and sentences $p_n$ such that
\[
 \ell_S(\Box_Sp_n\to p_n)=O(\log(n+2)),\qquad
 \ell_S(p_n)> n/C-1.
\]
Consequently, after taking a tail, this family satisfies the fixed-factor
inequality in MAIS-O11(1), including its additive toll, with $\delta=1$.
The main construction uses a countable, effectively coded set of new nullary
predicates and a primitive recursive verifier that performs very long finite
searches. An arithmetic-only variant is also given, with its failure of
uniform instantiation explicitly proved.
It addresses the stated permission to choose an efficient system, where
``efficient'' has the three-clause meaning in MAIS-A1. It does not establish
a speedup in the specified $\PAbin$, decide MAIS-O11(2), or give a
polynomial-time implementation of the specified verifier.
\end{abstract}

\paragraph{Scope and verification status.}
This is a mathematical construction, with its reductions and assumptions
spelled out below. Neither a byte-level implementation of its arithmetization
nor a Lean formalization has been produced. No claim of historical novelty or
independent expert verification is made. The deliberately slow checker is an
essential part of the example, rather than a performance detail that can be
discarded. The source question and its fixed-system polynomial-overhead
conjecture must be kept distinct.

\section{The precise target and the allowed freedom}

Write $\ell_T(A)$ for the least number of written characters in a $T$-proof
of $A$. All abbreviation definitions and names count. The requested lower
bound is a family $P_i$ with reflection lengths tending to infinity and
\begin{equation}\label{eq:target}
 \ell_S(P_i)\ge 2\ell_S(\Box_SP_i\to P_i)+C_1(|P_i|+1),
\end{equation}
where $C_1$ is a fixed constant for the cheap transformation from a proof of
$P$ to a proof of $\Box_SP\to P$. This gives $\delta=1$ in the problem.

MAIS-A1 calls a system efficient if it is a consistent recursively
axiomatized extension of PA, has logarithmic-length numerals, and allows
nested proof abbreviations. Its definition of a proof requires a primitive
recursive checker. These requirements do not impose a polynomial running-time
bound on the checker, nor do they require the language of an extension to
have the same finite signature as PA. Both freedoms are used here.

The proof files below use an ordinary Hilbert calculus. There is no L\"ob
inference rule built into the calculus and no extra padding charged to
conclusions. The additional computational work takes place when the verifier
checks whether the file is admissible.

\section{A fixed base calculus and its hard finite sentences}

Fix a Hilbert presentation $B$ of PA with efficient binary numerals and
nested abbreviations for complete terms and formulas. Choose a fixed finite
basis of propositional axiom schemes, the usual quantifier and equality
schemes, and PA induction, so checking an expanded proof takes polynomial
time in its expanded length. One may implement
binary constructors $b_0(t)=2t$ and $b_1(t)=2t+1$ as a fixed definitional
extension of arithmetic. Assume $B$ is consistent.

All tokens, including variable and abbreviation identifiers, are represented
over a fixed finite character alphabet. Identifiers use binary indices.
An abbreviation stands for a term or formula, not a fragment of a token;
in particular it cannot assemble a new predicate name from pieces of names.
Definitions can refer to previous definitions. A verifier may expand the
definitions and check the resulting Hilbert proof. This is primitive recursive,
and PA proves termination and the usual elementary syntactic invariants of
this specified procedure. Abbreviations need not be expanded when counting
the length of the input file.

Let $\Prf_B(u,v)$ be a natural, PA-provably $\Delta_1$ representation of this
checker, and let $\Len(u)$ count written characters. By the parameterized
diagonal lemma there is a fixed arithmetic formula $D(x)$ such that, writing
$D_n=D(\num n)$,
\begin{equation}\label{eq:finite-diagonal}
 B\vdash\forall x\bigl(D(x)\leftrightarrow H(x)\bigr),
\end{equation}
where
\begin{equation}\label{eq:H}
 H(n)\;:\Longleftrightarrow\;
 \text{no $B$-proof of $D_n$ has written length at most $n$}.
\end{equation}
Inside arithmetic the conclusion code in \eqref{eq:H} is produced by the
primitive recursive operation that substitutes the canonical numeral for $n$
into the fixed formula $D$. Thus \eqref{eq:finite-diagonal} is a uniform
arithmetical theorem, not an informal substitution of a free variable into a
quotation. The predicate $H$ is decided by checking all proof files of length
at most $n$; it has a PA-provably $\Delta_1$ representation. Also
$|D_n|=O(\log(n+2))$.

\begin{lemma}\label{lem:hard}
For every standard $n$, $H(n)$ is true, $B\vdash D_n$, and
\[
 \ell_B(D_n)>n.
\]
\end{lemma}
\begin{proof}
Suppose $u$ were a $B$-proof of $D_n$ of length at most $n$. Its existence
would give $B\vdash D_n$. Numeralwise verification of the particular finite
proof $u$ would also give $B\vdash\neg H(\num n)$. But
\eqref{eq:finite-diagonal} and the proof of $D_n$ give $B\vdash H(\num n)$.
This contradicts consistency. Hence no such $u$ exists and $H(n)$ is true.
The exhaustive finite check establishing $H(n)$ has a terminating computation;
PA can verify that computation, however long it is. Thus
$B\vdash H(\num n)$, and \eqref{eq:finite-diagonal} yields $B\vdash D_n$.
\end{proof}

This is the familiar finite G\"odel-sentence lower bound; see Pudl\'ak's
discussion of the diagonal sentences in \cite[Section 3.1]{pudlak2017}.
We have included its proof because it works for the chosen abbreviation
calculus without importing a quantitative finite-consistency theorem from
a different calculus. No bound on the time taken by the finite check is
used to infer that its proof is short.

\section{Constructing a verifier that refers to its own proof predicate}

Expand the language by distinct nullary predicates
\[
 p_0,p_1,p_2,\ldots.
\]
These are separate predicate symbols, not instances of a unary predicate
$p(x)$. Encode the name $p_n$ with a binary index, so $|p_n|=O(\log(n+2))$.
The characters in the name are charged, and names cannot be formed by
abbreviation expansion. Fix an injective primitive recursive formula coding
with
\begin{equation}\label{eq:atom-code}
 \#p_n=2n+1.
\end{equation}
Assign even codes to all other formulas using an ordinary string coding.
The numeral for a formula's code still has size $O(|A|+1)$. In particular,
the canonical numeral for $\#p_n$ is $b_1(\num n)$.
Proof-file codes can use the usual base-alphabet string encoding.

\subsection{The parameterized checker}

For a code $e$ of an arithmetic formula $\vartheta(u,v)$ with two designated
free variables, define $\Chk(e,\pi,y)$ by the following terminating
algorithm:
\begin{enumerate}
\item Parse $\pi$ as a Hilbert proof with the same abbreviation discipline as
$B$, and require its expanded conclusion to have code $y$.
\item Allow the ordinary logical and PA axioms in the expanded language,
modus ponens, generalization, and abbreviation definitions. In addition,
allow the axiom instances
\begin{equation}\label{eq:param-reflection}
 R_{e,n}:=\bigl(\exists u\,\vartheta(u,\num{2n+1})\bigr)\to p_n.
\end{equation}
The axiom formula must occur in the proof file, possibly through previously
defined abbreviations. Recognizing it only inspects the text coded by $e$;
it does not evaluate that formula.
\item For every primitive predicate $p_j$ occurring in the file, including
in abbreviation definitions, run the finite test $H(j)$ and require that it
returns true. An invalid formula code $e$ or malformed file is rejected.
\end{enumerate}

The algorithm is primitive recursive. The last step performs bounded proof
search in the already fixed base system $B$. There is no search for an
unbounded proof and no recursive invocation of the new verifier.

In $B$ one proves the syntactic invariant
\begin{equation}\label{eq:guard}
 B\vdash\forall e\,\forall\pi\,\forall n\,
 \bigl(\Chk(e,\pi,2n+1)\to H(n)\bigr).
\end{equation}
Indeed the only formula with code $2n+1$ is $p_n$. A file whose expanded
conclusion is $p_n$ must contain that literal predicate token in one of its
records or definitions: expansion of complete terms and formulas does not
manufacture new predicate tokens. Induction over the ordered abbreviation
definitions proves this fact in PA. The verifier has therefore checked
$H(n)$. This establishes \eqref{eq:guard} uniformly, without assuming that
every $H(n)$ is true.

\subsection{A fixed point for the checker description}

Apply the parameterized diagonal lemma to the arithmetic relation
$\Chk(e,\pi,y)$. Obtain an arithmetic formula $V(\pi,y)$, with code
$v=\#V$, for which
\begin{equation}\label{eq:verifier-fp}
 B\vdash\forall\pi\,\forall y\,
 \bigl(V(\pi,y)\leftrightarrow\Chk(\num v,\pi,y)\bigr).
\end{equation}
It can be chosen PA-provably $\Delta_1$: the diagonal construction uses
primitive recursive, PA-provably total substitution operations, and $\Chk$
is such a relation. Equivalently, substitute the graph of the unique output
of the diagonal substitution operation into its $\Delta_1$ definition.

Now define an $S$-proof to be a file accepted by the concrete checker
$\Chk(v,-,-)$. Define its arithmetized proof predicate to be $V$ and put
\[
 \Box_S A:=\exists\pi\,V(\pi,\gn A).
\]
Equation \eqref{eq:verifier-fp} identifies this predicate with the actual
checker inside PA. It is not merely an externally equivalent alternative
provability predicate. The extra axioms \eqref{eq:param-reflection} are now
exactly
\begin{equation}\label{eq:Ri}
 R_n=\Box_Sp_n\to p_n.
\end{equation}
This removes the apparent circularity in saying that the extra axioms
mention provability in the very system being defined.

\subsection{A finite runtime bound for the verifier}

The verifier is deliberately slow, but not an unspecified recursive search.
There is a fixed integer $c\ge1$ for which an implementation on proof files
of length $m$ computes the conclusion and acceptance result in at most
\[
 T(m)=2^{\,2^{c(m+1)^2}}
\]
steps. Comparing the conclusion against a separately supplied code also
charges the length of that code; alternatively reject an oversized supplied
code after reading a prefix longer than the maximal possible conclusion.

Here is a conservative bound sufficient for the claim. A file of length $r$
has at most $r$ ordered definitions. Each expansion stage increases the
current length by a factor at most $r+1$, so the expanded text has length at
most $(r+1)^{r+1}\le 2^{(r+1)^2}$ for $r\ge1$. There are at most
$|\Sigma_B|^{j+1}$ candidate proof files of length at most $j$. Their
expansion and polynomial-time checking therefore decide $H(j)$ within
$2^{c_2(j+1)^2}$ steps for a fixed $c_2$. In an $m$-character $S$-file,
every tested index $j$ has $O(m+1)$ binary digits, and there are at most $m$
distinct literal predicate names. Substitution into this bound, together with
the ordinary syntax checks, is absorbed by $T(m)$ after increasing $c$.
Malformed inputs can be rejected before the corresponding work is done.

All these bounds concern fixed finite loops, expansion, and arithmetic on
finite strings; PA proves termination with an elementary bound by the same
inductions. This is a runtime bound for the chosen implementation, not a
polynomial-time claim or an explicit length certificate for internalizing
the computation as an arithmetic proof.

Combining \eqref{eq:guard}, \eqref{eq:verifier-fp}, and
\eqref{eq:finite-diagonal} gives a \emph{single fixed $B$-proof} of
\begin{equation}\label{eq:uniform-image}
 \forall n\,\Bigl(\bigl(\exists\pi\,V(\pi,b_1(n))\bigr)\to D(n)\Bigr).
\end{equation}
This uniform theorem is the main reason the lower-bound interpretation below
can be efficient even though the verifier is not.

\section{Consistency, arithmetic strength, and L\"ob's theorem}

\begin{proposition}\label{prop:S}
The system $S$ is a consistent, recursively presented extension of $B$.
It is conservative over $B$ for arithmetic sentences, has logarithmic
numerals and nested priced abbreviations, and satisfies the
Hilbert--Bernays--L\"ob derivability conditions for $\Box_S$.
For each $n$, both $R_n$ and $p_n$ are $S$-theorems.
\end{proposition}
\begin{proof}
Every $B$-proof is an $S$-proof: it contains no new predicate symbols,
so the extra finite tests are vacuous. The new syntax and checker are
effective and primitive recursive. By Lemma~\ref{lem:hard}, all the tests
$H(j)$ are true for standard $j$. Thus each one-line extra-axiom proof
of $R_n$ is accepted.

For consistency, interpret every $p_j$ as the arithmetic tautology $0=0$,
leaving all arithmetic unchanged. Every extra axiom becomes
$\Box_Sp_j\to(0=0)$, a $B$-theorem. In this expression $\Box_Sp_j$ is
an arithmetic formula containing the numeral $\num{2j+1}$; its numeral
is left unchanged by the interpretation. Logical axioms, PA induction
instances, and inference rules translate to valid $B$-derivations.
Consequently an $S$-proof of an arithmetic sentence gives a $B$-proof of
that same sentence. In particular an $S$-contradiction would contradict
the consistency of $B$.

For completeness, the required derivability conditions hold for the
specified verifier, rather than being assumed for an arbitrary predicate.
\begin{enumerate}
\item If $S\vdash A$, the acceptance computation of its particular finite
proof can be verified in $B$. Hence $B\vdash\Box_S A$.
\item PA proves correctness of the primitive recursive proof-composition
operation for modus ponens. Combining two accepted proof files introduces
no new primitive predicate tokens; all the requisite tests $H(j)$ were
already checked in the input files. This proves
\[
 S\vdash\Box_S(A\to B')\to(\Box_S A\to\Box_S B').
\]
\item The formula $\Box_S A$ is arithmetical $\Sigma_1$, using a
$\Sigma_1$ presentation of $V$. Formalized $\Sigma_1$-completeness in PA
gives $\Box_S A\to\Box_B\Box_S A$. PA also proves that purely arithmetic
$B$-proofs embed into $S$, since the added tests are vacuous for those files.
It follows that $S\vdash\Box_S A\to\Box_S\Box_S A$.
\end{enumerate}
The diagonal lemma is available in the expanded language, as usual for an
extension of arithmetic. L\"ob's theorem therefore applies to $S$. Applying
it to its theorem $R_n$ proves $S\vdash p_n$.
\end{proof}

The tests do not add certificates or padding to the written proofs.
Combining and specializing proofs uses the underlying Hilbert syntax.
In particular, appending $P\to(\Box_S P\to P)$ and a modus-ponens line
gives the cheap-reflection inequality from MAIS-A1, with some fixed $C_1$.
The added formulas introduce no new $p_j$ outside those already in $P$:
the quotation in $\Box_S P$ is arithmetic text. One can omit explicit
line references in an MP record and let the checker find the premises,
so this argument does not assume that arbitrarily large line numbers
have constant character cost.

\section{The polynomial interpretation and the lower bound}

Define the interpretation $I$ by
\[
 I(p_j)=D_j,
\]
and let $I$ act identically on arithmetic and commute with logical
connectives and quantifiers. Every $D_j$ is closed. In particular, $I$
does \emph{not} replace a quoted old formula code by the code of its
interpretation. Therefore
\begin{equation}\label{eq:translated-reflection}
 I(R_j)=\bigl(\exists\pi\,V(\pi,\num{2j+1})\bigr)\to D(\num j),
\end{equation}
which is an instance of the uniform $B$-theorem
\eqref{eq:uniform-image}.

\begin{lemma}\label{lem:translate}
There is a fixed integer $C\ge1$ such that any $S$-proof $\pi$ of $A$,
of written length $m$, can be translated into a $B$-proof of $I(A)$
of written length at most $C(m+1)$.
\end{lemma}
\begin{proof}
Let $J$ be the finite set of indices of primitive predicates actually written
in $\pi$. Whole-expression abbreviations cannot create a new predicate token.
Consequently every predicate in an expanded record belongs to $J$, and
\[
 \sum_{j\in J}\bigl(1+\bits(j)\bigr)\le c_0(m+1)
\]
for a fixed encoding constant $c_0$. This charges each distinct index to one
of its literal occurrences in the input; it does not bound each index by $m$
and then unnecessarily pay that bound $m$ times.

For each $j\in J$, introduce a fresh $B$-abbreviation $d_j$ for
$D(\num j)$. Its name is a binary-indexed abbreviation name of length
$O(1+\bits(j))$, not a new primitive symbol. The total cost of these
definitions is $O(m+1)$. Translate the original definitions in order and the
original records by replacing $p_j$ with $d_j$ while retaining references to
earlier definitions. This costs $O(m+1)$ characters: each replacement name
has at most a constant multiple of the source predicate name's length.

Freshness can be implemented without a $\log m$ factor. Reserve disjoint
fixed residue classes of binary abbreviation indices for renamed source
names, the names $d_j$, and the finitely many names used by fixed proof
templates. Sending an index $i$ to $qi+r$, for fixed $q,r$, increases its
self-delimiting binary name length by at most a constant factor. No original
name is assumed to have constant size.

Before processing the translated input records, include one copy of the
fixed proof of \eqref{eq:uniform-image}. For each $j\in J$, instantiate
that theorem once at $\num j$ and retain the resulting proof of
\[
 I(R_j)=\bigl(\exists\pi\,V(\pi,\num{2j+1})\bigr)\to D(\num j).
\]
The canonical numeral for $2j+1$ is $b_1(\num j)$. Thus this preparation
costs $O(1+\bits(j))$ characters per index, and $O(m+1)$ altogether.
The preparation uses only the new definitions, so it can precede all source
definitions. It establishes each needed reflection image once, even if the
source proof invokes that axiom repeatedly.

Logical and arithmetic axioms translate to axioms; induction instances remain
induction instances. Modus ponens and generalization remain valid, and their
side conditions are preserved because the replacements $D_j$ are closed.
Replace an extra-axiom record by a reiteration of its already proved image.
No new rule is necessary: prove $F\to F$ by a fixed propositional template
and apply modus ponens to the earlier $F$. Here $F$ is written using the
translated representation of that source record, including any abbreviations
then in scope. The cost is a constant multiple of that record's written
length. Premise lookup uses equality of expanded formulas, so the cached
proof and this representation need not use the same abbreviation names.
We use the permitted version of Hilbert records in which the verifier finds
modus-ponens premises, rather than requiring explicit line references.

Summing these costs gives $C(m+1)$. The last translated inference concludes
$I(A)$, including when the source's last line is a reflection axiom. The
argument bounds written proof length without expanding the definition graph.
It does not assert a polynomial running time for the translator.
\end{proof}

\begin{theorem}\label{thm:main}
For the system $S$ above there are fixed constants $a,b,C\ge1$ such that,
for every $n\ge0$,
\begin{align}
 |p_n|+1&\le b(1+\bits(n)),\label{eq:pn-size}\\
 \ell_S(\Box_Sp_n\to p_n)&\le a(1+\bits(n)),\label{eq:reflection-upper}\\
 \ell_S(p_n)&>n/C-1,\label{eq:direct-lower}
\end{align}
where $\bits(n)$ is the binary digit length, with $\bits(0)=1$.
Moreover $\ell_S(\Box_Sp_n\to p_n)\to\infty$ as $n\to\infty$.
Thus a tail of this family satisfies \eqref{eq:target}.
\end{theorem}
\begin{proof}
The first bound follows from the binary coding of predicate names.
The second is the length of the one-line axiom proof of $R_n$.
The fixed arithmetic formula $V$ contributes a constant; the numeral
$\num{2n+1}$ and the name $p_n$ each contribute $O(\log(n+2))$.
The finite check required to accept this line is true by
Lemma~\ref{lem:hard}, and its computation time is not proof-file length.

Let $m=\ell_S(p_n)$, which is finite by Proposition~\ref{prop:S}.
Lemma~\ref{lem:translate} gives a $B$-proof of $D_n$ of length at most
$C(m+1)$. Lemma~\ref{lem:hard} implies
\[
 n<\ell_B(D_n)\le C(m+1),
\]
and hence \eqref{eq:direct-lower}.

For every fixed length bound $M$ there are only finitely many proof files
with at most $M$ characters. Each has one designated conclusion. The
sentences $R_n$ are pairwise syntactically distinct, so only finitely many
can have such a proof. This proves that their least proof lengths tend
to infinity, including in the presence of abbreviations.

Finally $n$ eventually dominates every fixed multiple of
$1+\bits(n)$. The lower bound therefore eventually exceeds
$2a(1+\bits(n))+C_1b(1+\bits(n))$, which implies
\eqref{eq:target}.
\end{proof}

\paragraph{An explicit rule selecting the witness sequence.}
The proof need not leave the choice of a tail informal. Fix integer constants
$a,b,C,C_1$ from the displayed proof templates and translation. Set
\[
 T(n)=(2a+C_1b)(1+\bits(n)).
\]
Let $n_0$ be the least positive integer satisfying
$n_0\ge C(T(n_0)+1)$, and let $n_{i+1}$ be the least integer greater
than $n_i$ satisfying the same decidable inequality. These searches terminate
because $n/(1+\bits(n))\to\infty$.
Put $P_i=p_{n_i}$. If an $S$-proof of $p_{n_i}$ had length at most $T(n_i)$,
its interpretation would contradict
$n_i<\ell_B(D_{n_i})$. Thus
\[
 \ell_S(P_i)>T(n_i)
 \ge 2\ell_S(\Box_SP_i\to P_i)+C_1(|P_i|+1).
\]
This specifies an effective family and the concrete factor $\delta=1$;
numerical values of the constants await an implementation of the fixed
proof templates, rather than being asserted without a character ledger.

\section{Adversarial checks and the limits of the result}

\paragraph{No axiom is accepted by an oracle.}
The only unusual acceptance condition is the finite computation $H(j)$.
It halts even if the consistency assumption is dropped. Consistency is used
to show that its standard output is true, not to make it computable.

\paragraph{The proof predicate is the right one.}
Equation \eqref{eq:verifier-fp} is a PA theorem identifying $V$ with the
actual parameterized verifier at its fixed description. The reflection
axioms use $\Box_S$, not provability in a stronger background theory or
an unaligned surrogate. The derivability conditions were checked separately.

\paragraph{The use of L\"ob is legitimate.}
L\"ob is invoked after consistency, the actual accepted reflection proofs,
and the derivability conditions have been established. It is not used to
justify acceptance of an axiom or the finite G\"odel lower bound.

\paragraph{The lower bound quantifies over every direct proof.}
The interpretation applies to an arbitrary $S$-proof, including one using
other reflection axioms, induction, or an unexpected shortcut. Its output
is a forbidden short proof of $D_n$ if the input is too short. Merely comparing
two particular proof scripts would not establish this.

\paragraph{Abbreviations have not been silently expanded.}
The translation preserves the definition graph. Closed replacements preserve
binding conditions. Identifier characters are charged. Restricting
abbreviations to complete expressions is a declared choice of system and
prevents a short file from manufacturing a primitive predicate name of
exponentially many digits.

\paragraph{The bottleneck is internalizing an accepted proof.}
A one-line proof of $R_n$ is accepted after the external checker completes
$H(n)$. To prove internally that this particular file is accepted, arithmetic
must establish the relevant check. The existence of a primitive recursive
checker gives no polynomial bound on the length of that internal
verification. The example exploits precisely this missing requirement.

Although every standard test $H(n)$ succeeds, $B$ cannot prove
$\forall n\,H(n)$. Otherwise \eqref{eq:finite-diagonal} would give a fixed
proof of $\forall n\,D(n)$, whose numeral instances would have
$O(\log(n+2))$-length proofs, contrary to Lemma~\ref{lem:hard}.
Consequently one cannot simply erase the guard inside arithmetic because
it is extensionally redundant. The choice of the actual verifier and its
arithmetization is essential.

This is not a lower bound on the complexity of the \emph{set} of accepted
files. An externally equivalent algorithm can omit all the true tests and
check only the underlying Hilbert derivation. Its arithmetic acceptance
formula is then different. An external equality of accepted files does not
supply a PA proof identifying that formula with $V$. Thus a slow specified
algorithm, a polynomial-time decidable proof relation, and an efficiently
internalizable provability formula are three different assertions.

\paragraph{Why this is not a solution for the default PA presentation.}
The system has new primitive sentence symbols, extra reflection axioms, and
a slow verification condition. The ability to abbreviate expressions, by
itself, does not license replacing the specified $\PAbin$ by this system
when answering a question that fixes $\PAbin$. Theorem~\ref{thm:main}
answers only the explicitly allowed chosen-system branch of MAIS-O11(1).

\paragraph{What follows for the overhead function.}
For $k_n=a(1+\bits(n))$ and $r_n=b(1+\bits(n))$, the definition in
MAIS-A1 gives
\[
 F_S(k_n,r_n)\ge\ell_S(p_n)>n/C-1.
\]
Thus this particular $F_S$ is not bounded by any polynomial in its two
arguments. This shows that the three-clause definition of efficiency alone
cannot force polynomial L\"ob overhead. It does \emph{not} decide
the conjecture for $F_{\PAbin}$ in MAIS-O11(2)/MAIS-O10.

\paragraph{What is not being claimed about the companion problems.}
The elementary runtime bound above is not a bound on the shortest arithmetic
proof verifying acceptance. No tight expansion bound, internal budget certificate with explicit constants,
FairBot threshold, or bounded-$\mathsf{GL}$ theorem is obtained here.
The resource-bounded companion statements have further uniform certificate
hypotheses. This note establishes the unbounded chosen-system separation;
it does not substitute that result for those additional obligations.
The chosen verification procedure is slow; no lower bound against all
extensionally equivalent checking algorithms is asserted. Replacing its
arithmetization by the one associated with a different checker changes the
reflection formula under study. The arithmetic-only variant avoids the
language extension, but fails uniform instantiation. No answer for the
standard, fixed-verifier problem is asserted.

\section{An arithmetic-only variant and an exact obstruction}
\label{sec:arithmetic-variant}

The new predicate symbols can be avoided if one asks only for the unbounded
L\"ob conditions and the three-clause definition of efficiency. This variant
has exactly the same accepted proof files as $B$, under the consistency
assumption. It nevertheless uses a different checker and a different
arithmetization of provability. It fails the agenda's additional uniform
instantiation requirement. Both facts are essential to its interpretation.

\subsection{A recognizable hard arithmetic family}

Reserve the arithmetic variable $z_0$. A variant of the first diagonal
construction gives a fixed formula
\[
 G(x)=\forall z_0\,\bigl(x=x\wedge E(x)\bigr)
\]
with $z_0$ absent from $E$, and a PA-provably $\Delta_1$ predicate $K$, such
that
\begin{equation}\label{eq:arithmetic-diagonal}
 B\vdash\forall x\,(G(x)\leftrightarrow K(x)),\qquad
 K(n)\equiv\text{no $B$-proof of $G(\num n)$ has length at most $n$}.
\end{equation}
For clarity, this is not an arbitrary syntactic modification of an existing
fixed point. Apply the parameterized diagonal lemma to the operator that,
given a formula code $e$, first wraps that formula as
$\forall z_0(x=x\wedge\cdot)$ and then tests for short proofs of its numeral
instances. The resulting $E$ asserts the absence of a short proof of the
wrapped formula $G$. Since the wrapper is provably equivalent to $E$, this
gives \eqref{eq:arithmetic-diagonal}. The diagonal construction can use
variables other than $z_0$. The visible conjunct $x=x$ ensures that different
canonical numeral instances $G_n:=G(\num n)$ are syntactically distinct.

Exactly as in Lemma~\ref{lem:hard}, consistency gives, for every standard $n$,
\begin{equation}\label{eq:arithmetic-hard}
 K(n),\qquad B\vdash G_n,\qquad \ell_B(G_n)>n.
\end{equation}
In particular $B\nvdash\forall n\,K(n)$: such a theorem, combined with
\eqref{eq:arithmetic-diagonal}, would give $O(\log(n+2))$-length proofs of
all the $G_n$ by numeral instantiation, contradicting the last inequality.

Use a new formula coding $c_*$ in which $c_*(G_n)=2n+1$ and every other
formula has an even ordinary string code. This coding is injective, primitive
recursive, and has linear-size quoted numerals. The base coding used inside
$K$ is unchanged. This separation of the two codings avoids a circular
definition of $G$.

\subsection{The actual arithmetic checker}

Define $\Prf_*(\pi,y)$ to mean the following. Decode $y$ using $c_*$; check
that $\pi$ is a $B$-proof of that formula; expand its definitions; and, for
every subformula of every expanded record or definition that is exactly a
canonical instance $G_n$, run the finite test $K(n)$ and require success.
This is a primitive recursive, PA-provably terminating algorithm. All these
tests succeed on standard inputs by \eqref{eq:arithmetic-hard}. Thus, if
$S_*$ is the system checked by this algorithm, then
\begin{equation}\label{eq:identical-files}
 \{\text{$S_*$-proof files}\}=\{\text{$B$-proof files}\},\qquad
 \ell_{S_*}(A)=\ell_B(A).
\end{equation}
The equality is an external theorem from consistency. It is not an
equivalence of the two proof predicates provable in $B$.

Choose a $\Sigma_1$ presentation $V_*(u,y)$ of this PA-provably $\Delta_1$
checker, renaming all its variables to avoid $z_0$, and set
$\Box_*A:=\exists u\,V_*(u,\num{c_*(A)})$. The existential binder also
avoids $z_0$. No quoted formula occurs literally inside its quotation, so
\emph{every} formula $\Box_*A$ contains no variable token $z_0$.

The guard gives a fixed uniform theorem
\begin{equation}\label{eq:arithmetic-reflection}
 B\vdash\forall n\,\Bigl((\exists u\,V_*(u,b_1(n)))\to G(n)\Bigr).
\end{equation}
Indeed an accepted proof with conclusion code $2n+1$ has conclusion $G_n$,
so its conclusion itself triggers the test $K(n)$. Use
\eqref{eq:arithmetic-diagonal} to finish. Instantiating the fixed proof at
$\num n$ costs $O(\log(n+2))$ characters. By
\eqref{eq:identical-files}, those are also $S_*$-proofs. Hence
\begin{equation}\label{eq:arithmetic-separation}
 \ell_{S_*}(\Box_*G_n\to G_n)=O(\log(n+2)),\qquad
 \ell_{S_*}(G_n)>n.
\end{equation}
Reflection lengths tend to infinity by the same finite-file argument as
before. Appending a cheap reflection proof introduces no new guarded
subformula: $\Box_*A$ has no $z_0$, and all guarded subformulas in $A$ were
already checked in the given proof. Thus the additive toll also remains
available, and a tail of this family satisfies the requested $\delta=1$
inequality.

\subsection{Why L\"ob's theorem still applies}

This is not merely a predicate chosen to make a reflection statement easy.
The following verifies all three unbounded derivability conditions.
\begin{enumerate}
\item An actual accepted proof has a finite acceptance computation verifiable
in $B$. By \eqref{eq:identical-files}, the resulting theorem is also an
$S_*$-theorem. This is external necessitation.
\item Concatenate accepted proofs of $A\to F$ and $A$, rename abbreviation
identifiers to avoid collisions, and append $F$ by modus ponens. The appended
formula $F$ is already a subformula of the first input's conclusion. Thus no
new $G_n$ test is required. PA proves this syntactic fact and consequently
proves the distribution axiom for $\Box_*$. The proof itself belongs to
$S_*$ by \eqref{eq:identical-files}.
\item Formalized $\Sigma_1$-completeness for the ordinary base checker gives
\[
 B\vdash\Box_*A\to\Box_B(\Box_*A).
\]
It remains to embed these particular $B$-proofs into $S_*$, without assuming
a general proof-predicate equivalence. Fix a finite set of variable names
containing every variable in the template $\exists u\,V_*(u,y)$ and excluding
$z_0$. There is a primitive recursive injective renaming of all arithmetic
variables which fixes this finite set and has $z_0$ outside its image. For
example, fix indices $1,\ldots,M$, send $0$ to $M+1$, and send $i>M$ to
$i+1$. Apply this renaming throughout a $B$-proof of $\Box_*A$ and its
definitions. Logical axioms, PA induction, and inference rules remain valid;
the conclusion is unchanged. The renamed file has no $z_0$, so it has no
subformula $G_n$ and all additional tests are vacuous. PA proves correctness
of this transformation. Therefore
\[
 B\vdash\Box_B(\Box_*A)\to\Box_*(\Box_*A),
\]
and hence $S_*\vdash\Box_*A\to\Box_*\Box_*A$.
\end{enumerate}
The diagonal lemma is available for the new effective coding. These facts
establish applicability of ordinary L\"ob's theorem. The proof of the
separation itself already supplies theorems on both sides independently.

\subsection{The exact uniform property that fails}

Let $\tau(x):=G(x)\to G(x)$. There is a fixed $S_*$-proof of
$\forall x\,\tau(x)$, since it is a $B$-theorem and the accepted proof files
coincide. But
\begin{equation}\label{eq:uniform-failure}
 S_*\nvdash\forall n\,\Box_*\tau(\num n).
\end{equation}
The expression here denotes the usual arithmetized numeral-instance coding.
To prove the claim, an accepted proof of $\tau(\num n)$ contains $G_n$ as a
subformula of its conclusion. PA therefore proves uniformly that
$\Box_*\tau(\num n)\to K(n)$. A theorem excluded by
\eqref{eq:uniform-failure} would give a $B$-proof of $\forall n\,K(n)$,
contradicting \eqref{eq:arithmetic-hard}.

Thus even the unbounded conclusion of the agenda's quantifier-distribution
principle fails here; its bounded version cannot hold. This variant answers
only the literal unbounded, chosen-presentation question. It must not be used
to instantiate the stronger resource-bounded hypotheses. The main
construction avoids this particular obstruction: substituting a numeral for
an arithmetic variable does not create a new nullary predicate symbol.

The possibility that externally identical theories have internally different
provability predicates is familiar from the slow-provability literature; see
Henk and Pakhomov \cite{henk-pakhomov}. That literature does not, by itself,
establish the particular size separation above. The finite tests, syntactic
renaming argument, and lower bound are the required additional ingredients.

\section{Rejected mechanisms and a precise formalization target}

Two simpler constructions do not answer the intended question. First,
requiring extra padding only for direct conclusions creates a formal size
gap but can destroy cheap proof composition. The present construction
adds no such padding. Second, a finite G\"odel sentence by itself supplies
a lower bound on $\ell_B(D_n)$, but supplies no short proof of
$\Box_BD_n\to D_n$. The new verifier and interpretation are what bridge
that missing step in the chosen system.

Another attempted route is to make a sentence assert that every proof of
it has a sufficiently shorter reflection proof. Consistency may then enforce
the desired comparison \emph{if the sentence is provable}, but it does not
establish that provability. Unprovability of both sentences is compatible
with the proposed comparison. The present construction supplies actual
reflection proofs and checks L\"ob's applicability instead.

For a Lean development, the following are separate mathematical obligations:
\begin{enumerate}
\item Define the base and extended syntax, proof files, expansion, actual
character lengths, and the primitive recursive checkers.
\item Formalize the uniform finite diagonal family, prove the hard-sentence
lemma from consistency, and formalize the second diagonal construction for
$V$. These require syntax and arithmetization, not merely a modal box axiom.
\item Prove the guard invariant, consistency/conservativity, and the
derivability conditions for the actual checker.
\item Implement the interpretation on compressed proof files and verify its
length bound. Prove that it handles every admitted rule and side condition.
\item Derive the explicit family inequality and audit all assumptions of the
final theorem. Consistency of the chosen PA presentation must be supplied
or proved in the metatheory; it must not disappear behind an undeclared axiom.
\end{enumerate}
Formalizing only the last inequality while assuming the guard and translation
lemmas would certify a conditional transfer theorem, not the whole construction.
This note contains no Lean files and no claim that Lean has checked it.

\begingroup\small
\begin{thebibliography}{9}
\bibitem{agenda}
\emph{Quantitative bounded L\"ob: proof-length overheads in L\"ob's theorem
and L\"obian cooperation}, MAIS-A1, uploaded version; and
\emph{Does L\"ob's rule shorten proofs?}, MAIS-O11, uploaded version.
The chosen-system permission is in Question 3.6(1) / MAIS-O11(1).

\bibitem{critch}
A. Critch, \emph{A parametric, resource-bounded generalization of L\"ob's
theorem, and a robust cooperation criterion for open-source game theory},
Journal of Symbolic Logic 84 (2019), 1368--1381.
\url{https://arxiv.org/abs/1602.04184}.

\bibitem{pudlak2017}
P. Pudl\'ak, \emph{Incompleteness in the finite domain},
Bulletin of Symbolic Logic 23 (2017), 405--441.
\url{https://arxiv.org/abs/1601.01487}.

\bibitem{pudlak1998}
P. Pudl\'ak, \emph{The lengths of proofs}, in S. R. Buss (ed.),
Handbook of Proof Theory, Elsevier, 1998, 547--637.
Section 2.1 uses polynomial-time decidability of the proof relation.
That extensional complexity property must be distinguished from the choice
of an arithmetic formula representing the relation.
\url{https://users.math.cas.cz/~pudlak/length.pdf}.

\bibitem{henk-pakhomov}
P. Henk and F. Pakhomov, \emph{Slow and Ordinary Provability for Peano
Arithmetic}, 2016 preprint, version 2.
\url{https://arxiv.org/abs/1602.01822}.
\end{thebibliography}
\endgroup
\end{document}
